% !TEX root = ../main.tex
\chapter{Kết luận và Hướng phát triển}
\label{Chapter5}

\section{Tóm tắt đóng góp}
\label{sec:tom-tat-dong-gop}

Luận văn đặt ra câu hỏi: liệu cầu nối Q-Former -- với cơ chế cross-attention có kiểm
soát -- có thể căn chỉnh tín hiệu cộng tác vào không gian ngữ nghĩa của LLM hiệu quả
hơn lớp chiếu tuyến tính của CoLLM hay không? Để trả lời, luận văn thiết kế, hiện thực và đánh
giá mô hình CoQLLM, đóng góp theo hai hướng đã phát biểu ở Mục~\ref{sec:dong-gop}:

\textbf{Đóng góp 1: Cầu nối Q-Former được thích nghi và tăng cường cho bài toán gợi ý.}
CoQLLM thay thế lớp MLP một lớp của CoLLM bằng kiến trúc Q-Former, trong đó tín hiệu cộng tác
(nhúng người dùng, sản phẩm mục tiêu, lịch sử) thay thế đặc trưng ảnh làm đầu vào
cross-attention và được đưa vào LLM qua đường soft token. Trên nền đó, luận văn bổ sung hai
đòn bẩy huấn luyện -- điều kiện hóa truy vấn theo người dùng và mục tiêu bảo toàn thứ tự
trong từng người dùng -- cùng các lựa chọn thiết kế (biểu diễn đa token, backbone đã tinh chỉnh chỉ
dẫn, tách lớp chiếu người dùng/sản phẩm). Kết quả chi tiết được trình bày ở
Chương~\ref{Chapter4}.

\textbf{Đóng góp 2: Đánh giá có hệ thống dưới quy trình AUC/uAUC trên hai tập dữ liệu.}
Luận văn đánh giá cầu nối Q-Former dưới quy trình CTR với AUC và uAUC -- cùng quy trình với
CoLLM -- trên MovieLens-1M và Amazon-Book, cho phép so sánh trực tiếp với CoLLM, BinLLM và
CoRA. Phân tích ablation tách biệt đóng góp của từng đòn bẩy, và phân rã warm/cold làm rõ
trần uAUC nằm ở chất lượng nhúng MF.

\section{Kết quả đạt được}
\label{sec:ket-qua-dat-duoc}

Từ kết quả ở Bảng~\ref{tab:ket-qua-chinh}, Bảng~\ref{tab:ket-qua-book} và
Bảng~\ref{tab:ablation}, có thể rút ra các kết luận sau:

\begin{itemize}
    \item \textbf{MovieLens-1M -- vượt SOTA trên cả hai độ đo.} Cấu hình đề xuất
    (+ \texttt{user\_conditioned}) đạt AUC $0{,}7475$ và uAUC $0{,}6968$, vượt BinLLM
    ($0{,}7425$/$0{,}6956$) trên cả AUC lẫn uAUC và vượt rõ CoLLM. Đây là điểm đáng chú ý
    vì nhiều phương pháp cùng hướng chỉ thắng một trong hai độ đo.

    \item \textbf{Amazon-Book -- cạnh tranh/vượt baseline trên miền khó.} Với
    \texttt{user\_conditioned} (Run-1), CoQLLM đạt $0{,}8147$/$0{,}5958$, vượt CoLLM-MF và
    BinLLM (theo lần chạy lại của CoRA) trên cả hai độ đo. Bổ sung \texttt{align\_rank\_loss}
    (Run-2) nâng uAUC lên $0{,}6062$, tiệm cận CoRA-MF ($0{,}6262$).

    \item \textbf{Hai đòn bẩy tác động lên hai độ đo khác nhau.} \texttt{user\_conditioned}
    là đòn bẩy AUC (biến thiên giữa các người dùng); \texttt{align\_rank\_loss} là đòn bẩy
    uAUC (bảo toàn thứ tự trong người dùng). Việc bật theo từng tập dữ liệu là siêu tham số
    hợp lệ theo miền.

    \item \textbf{Trần uAUC nằm ở MF teacher, không phải cầu nối.} Phân rã warm/cold cho
    thấy uAUC nhóm warm cao hơn hẳn nhóm cold trên cả hai tập dữ liệu (ví dụ Amazon-Book:
    $0{,}6238$ vs $0{,}5258$), cho thấy đòn bẩy cải thiện tiếp theo nằm ở chất lượng véc-tơ cộng tác
    cho người dùng/sản phẩm nguội (cold).
\end{itemize}

\section{Hạn chế}
\label{sec:han-che-ch5}

Luận văn nêu rõ các hạn chế sau một cách minh bạch:

\begin{enumerate}
    \item \textbf{Thiếu khoảng tin cậy:} Kết quả từ một seed duy nhất. Biến động thống
    kê giữa các lần chạy chưa được đo.

    \item \textbf{Thiếu ablation attention-MLP:} Chưa có hàng ablation thay Q-Former bằng
    bộ chiếu attention-MLP (không truy vấn học được) -- đối chứng quan trọng nhất để cô lập
    đóng góp của cấu trúc Q-Former. Cùng với các ablation thành phần khác (đa token, tách
    projector, \texttt{align\_rank} trên MovieLens-1M), đây là phần thực nghiệm còn để ngỏ.

    \item \textbf{Chi phí tính toán cao hơn CoLLM:} Cầu nối Q-Former làm tăng chi phí huấn
    luyện so với lớp MLP đơn giản của CoLLM (dù không sinh trọng số theo từng mẫu như CoRA).

    \item \textbf{Baseline không đồng nhất:} SellaRec dùng backbone LLM khác họ, và trên
    Amazon-Book số liệu baseline lấy từ lần chạy lại của CoRA (chọn checkpoint theo AUC,
    khác CoQLLM chọn theo uAUC) -- giới hạn khả năng so sánh thuần về cơ chế căn chỉnh.

    \item \textbf{Trần uAUC ở nhóm cold:} uAUC bị chặn trên bởi chất lượng nhúng MF cho
    người dùng/sản phẩm nguội, không phải bởi cầu nối.
\end{enumerate}

\section{Hướng phát triển}
\label{sec:huong-phat-trien}

Từ các hạn chế và quan sát trên, luận văn đề xuất những hướng nghiên cứu tiếp theo:

\textbf{Mở rộng thực nghiệm:}
\begin{itemize}
    \item Đánh giá trên Amazon Books và các tập dữ liệu ngoài miền phim để kiểm chứng
    khả năng tổng quát.
    \item Chạy nhiều seed và báo cáo mean $\pm$ std để đánh giá độ ổn định.
\end{itemize}

\textbf{Ablation và phân tích sâu hơn:}
\begin{itemize}
    \item Bổ sung hàng ablation attention-MLP (bộ chiếu có self-attention nhưng không có
    truy vấn học được) để cô lập đóng góp của cấu trúc Q-Former khỏi phần phi tuyến +
    attention đơn thuần.
    \item Phân tích attention map của Q-Former để hiểu truy vấn nào học cách trích xuất
    thông tin gì từ tín hiệu cộng tác.
\end{itemize}

\textbf{Cải thiện uAUC:}
\begin{itemize}
    \item Tiền huấn luyện MF trên dữ liệu dày đặc hơn (nhiều tháng hơn), đặc biệt cho
    nhóm người dùng/sản phẩm nguội -- đòn bẩy được phân rã warm/cold chỉ ra là trần thực sự
    của uAUC.
    \item Mở rộng mục tiêu bảo toàn thứ tự (\texttt{align\_rank\_loss}, đã hiệu quả trên
    Amazon-Book) sang các miền khác và khảo sát trọng số $\lambda$ tối ưu theo từng tập.
\end{itemize}

\textbf{Mở rộng kiến trúc:}
\begin{itemize}
    \item Bổ sung hàng ablation attention-MLP và các ablation thành phần còn để ngỏ dưới
    cùng backbone Vicuna để cô lập đầy đủ đóng góp của cấu trúc Q-Former.
    \item Khảo sát hướng sinh trọng số kiểu CoRA như một đường \textit{bổ sung} cho soft
    token (chưa hiện thực trong luận văn này) để so sánh trực tiếp hai cơ chế đưa tín hiệu cộng tác vào LLM.
    \item Khảo sát các backbone LLM khác (Llama-3, Mistral) để kiểm chứng tính tổng
    quát của cầu nối Q-Former.
\end{itemize}
